\begingroup
\raggedbottom
\setlength{\emergencystretch}{2em}
\section{概形的推出}\label{src:B.4}
\newterm{推出}{pushout} 是纤维积的对偶概念. 与纤维积不同, 推出并不总是存在. 不过, Ferrand 找到了一种保证推出存在的一般情形: 一个映射为闭浸入, 另一个映射仿射. Artin 环的 Ferrand 推出在形变理论中尤为重要, 例如 Rim--Schlessinger 判据 (\bookref{C.4.6}) 和 Artin 代数性公理 (\bookref{C.7.4}) 中出现的齐性条件. 我们还将使用推出来构造稳定曲线模空间之间的粘合态射 (\bookref{6.6.16}).

\subsection{推出的存在性}\label{src:B.4.1-subsection}
\begin{theorem}[Ferrand 推出]\label{src:B.4.1}
考虑概形图
\begin{equation}\label{src:B.4.2}
\begin{tikzcd}[column sep=3.5em,row sep=2.2em]
X_0\arrow[r,hook,"i"]\arrow[d,"f_0"']\arrow[dr,phantom,"\lrcorner" description]&X\arrow[d,"f"]\\
Y_0\arrow[r,hook,"j"']&Y
\end{tikzcd}
\end{equation}
其中 $i:X_0\hookrightarrow X$ 为闭浸入, $f_0:X_0\to Y_0$ 仿射. 若满足
\begin{quote}
\hypertarget{pushout-affine-neighborhood}{($\star$)} 对每个点 $y_0\in Y_0$, 都存在仿射开子概形 $U\subseteq X$, 使得 $f_0^{-1}(y)\subseteq U\cap X_0$,
\end{quote}
则存在概形 $Y$, 闭浸入 $j:Y_0\hookrightarrow Y$ 及仿射概形态射 $f:X\to Y$, 使得 \bookeqref{B.4.2} 为推出. 此外,
\begin{enumerate}[(a)]
\item 方块 \bookeqref{B.4.2} 为 Cartesian 图, $X\to Y$ 限制为同构 $X\setminus X_0\to Y\setminus Y_0$, 且拓扑空间 $|Y|$ 等同于推出 $|X|\amalg_{|X_0|}|Y_0|$;
\item 诱导映射
\[
\mathcal O_Y\longrightarrow
j_*\mathcal O_{Y_0}\times_{(j\circ f_0)_*\mathcal O_{X_0}}f_*\mathcal O_X
\]
为层的同构;
\item 若 $f_0$ 有限, 则 $f$ 也有限. 此时, 若 $X_0$, $X$ 和 $Y_0$ 均在某个 Noether 概形上局部有限型, 则 $Y$ 也如此.
\end{enumerate}
\end{theorem}
\begin{proof}
见 \cite[定理 $5.4$ 及 $7.1$]{Fer03}, \cite[定理 $6.1$]{Art70} 及 \cite[Tag 0ECH]{SP}.
\end{proof}
我们称 $Y=X\amalg_{X_0}Y_0$ 为\newterm{Ferrand 推出}{Ferrand pushout}. 注意, 若存在形如 \bookeqref{B.4.2} 的概形 Cartesian 图, 且 $f:X\to Y$ 仿射, 则条件 \hyperlink{pushout-affine-neighborhood}{($\star$)} 成立.
\begin{remark}[一般情形下推出的存在性]\label{src:B.4.3}
条件 \hyperlink{pushout-affine-neighborhood}{($\star$)} 并不总成立. 例如, 考虑 $f_0:X_0\to\Spec k$ 和 $i:X_0\hookrightarrow X$, 其中 $X$ 为代数闭域 $k$ 上的光滑固有 (但非射影) 三维簇, 且存在由两个 $k$-点组成的集合 $X_0$, 它们不包含在同一个仿射开集中. 不过, 推出作为代数空间总是存在, 且是代数空间范畴中的推出.
\end{remark}
\begin{example}[仿射情形]\label{src:B.4.4}
在仿射情形 $X=\Spec A$, $X_0=\Spec A_0$, $Y_0=\Spec B_0$ 下, $\Spec(A\times_{A_0}B_0)$ 就是推出 $X\amalg_{X_0}Y_0$.
\end{example}
\begin{example}[粘合与捏合]\label{src:B.4.5}
若 $X_0\hookrightarrow X$ 和 $X_0\hookrightarrow Y_0$ 均为闭浸入, 则推出 $X\amalg_{X_0}Y_0$ 可以视为 $X$ 与 $Y_0$ 沿 $X_0$ 的粘合. 例如, 结点曲线 $\Spec k[x,y]/xy$ 是两条 $\mathbb A^1$ 沿各自原点的并. 若 $X_0=Z\amalg Z$ 是 $X$ 中两个同构且不交的子概形的并, $X_0\to Z$ 为投影, 则推出 $X\amalg_{Z\amalg Z}Z$ 可以视为将 $X$ 中的两个 $Z$ 的副本\newterm{捏合}{pinching} 在一起. 例如, 结点三次曲线由捏合 $\mathbb P^1$ 中的 $0$ 和 $\infty$ 得到.
\end{example}
\begin{exercise}\label{src:B.4.6}
若 $X$ 为两个闭子概形 $Z_1$ 和 $Z_2$ 的概形论并, 证明 $X=Z_1\amalg_{Z_1\cap Z_2}Z_2$.
\end{exercise}
\newpage
\begin{example}[非 Noether 性]\label{src:B.4.7}
当 $f_0:X_0\to Y_0$ 仿射但不有限时, 推出 $X\amalg_{X_0}Y_0$ 往往不是 Noether 概形. 例如, 若 $X_0=V(x)\subseteq X=\mathbb A^2_k$, 且 $f_0:X_0\to\Spec k$, 则推出是由下列环定义的非 Noether 仿射概形:
\[
k[x,y]\times_{k[x]}k=k[x,xy,xy^2,xy^3,\ldots]\subseteq k[x,y].
\]
另一方面, 我们本就不应期望推出有限型: $\mathbb A^2_k$ 中的 $y$ 轴不能被收缩.
\end{example}

\subsection{推出的性质}\label{src:B.4.2-subsection}
\begin{proposition}[推出的性质]\label{src:B.4.8}
设 $X_0\hookrightarrow X$ 为闭浸入, $X_0\to Y_0$ 为仿射概形态射.
\begin{enumerate}[(1)]
\item 若 $Y\simeq X\amalg_{X_0}Y_0$ 为概形的 Ferrand 推出, 则自然函子
\[
\QCoh(Y)\longrightarrow\QCoh(Y_0)\times_{\QCoh(X_0)}\QCoh(X)
\]
限制到平坦 $\mathcal O$-模 (相应地, 有限型平坦 $\mathcal O$-模, 有限呈示平坦 $\mathcal O$-模) 构成的满子范畴上时为等价.
\end{enumerate}
考虑概形的交换立方体
\[
\begin{tikzcd}[column sep=1.8em,row sep=1.5em]
&X'_0\arrow[rr,hook]\arrow[dl]\arrow[dd,dashed]&&X'\arrow[dl]\arrow[dd]\\
X_0\arrow[rr,hook]\arrow[dd]&&X\arrow[dd]&\\
&Y'_0\arrow[rr,hook,crossing over]\arrow[dl]&&Y'\arrow[dl]\\
Y_0\arrow[rr,hook]&&Y&
\end{tikzcd}
\]
其中 $X_0\hookrightarrow X$ 为闭浸入, $X_0\to Y_0$ 仿射.
\begin{enumerate}[(1)]
\setcounter{enumi}{1}
\item 假设 $Y'\to Y$ 平坦, 且 $X'_0$, $Y'_0$ 和 $X'$ 均为基变换. 若 $Y\simeq X\amalg_{X_0}Y_0$, 则 $Y'\simeq X'\amalg_{X'_0}Y'_0$. 若 $Y'\to Y$ 为 fppf 态射, 则逆命题成立.
\item 若顶面与左侧面为 Cartesian 图, 且前面与后面均为 Ferrand 推出, 则所有面均为 Cartesian 图. 此外, 若 $Y'_0\to Y_0$ 和 $X'\to X$ 平展 (相应地, 光滑), 则 $Y'\to Y$ 也如此.
\item 假设 $Y$ 定义在概形 $S$ 上. 设 $S'\to S$ 为概形态射, $Y'$, $X'_0$, $Y'_0$ 和 $X'$ 为基变换. 若 $X_0\to S$ 平坦且 $Y\simeq X\amalg_{X_0}Y_0$, 则 $Y'\simeq X'\amalg_{X'_0}Y'_0$.
\end{enumerate}
\end{proposition}
\begin{proof}
(1)--(3) 由 \cite[定理 $2.2$]{Fer03} 推出; 另见 \cite[Tag 0D2K]{SP} 及 \cite[\S$4$]{AHHLR24}. (4) 是初等的: 化归到仿射情形 $Y=\Spec B_0\times_{A_0}A$ 和 $S=\Spec R$. 由于 $A_0$ 在 $R$ 上平坦, $R$-模的正合列
\[
0\longrightarrow B\longrightarrow A\times B_0\longrightarrow A_0\longrightarrow0
\]
与一个 $R$-代数取张量积后仍正合.
\end{proof}
\par\clearpage
\endgroup
